\section{国际贸易术语解释通则}
\begin{enumerate}
  \item 基本内容、价格构成、风险转移、交货地点、运输、投保、清关手续。
  \item 2020 年通则。11 个贸易术语。
        \begin{enumerate}
          \item 将与之相关的安保要求明确规定了，将安保费用纳入到运输费用中。
          \item 有义务承担运输的一方，可以自运 \hz 第三方承运。
          \item 掐头去尾。
                \begin{enumerate}
                  \item EXW(EX works, 工厂交货)。卖方义务最小，出门不管。
                  \item DDP(Delivered Duty Paid, 完税后交货)。卖方义务最大，送货上门。
                  \item 其余 9 个术语卖方办理出口手续，买方办理进口手续。\al{上海甲公司与 M 国乙公司签订 CFR 两用物项出口合同}。
                \end{enumerate}
          \item 最最最常用的贸易术语到目前为止 3 个。FOB、CFR、CIF。与船运（海运/河运）密切。
                \begin{enumerate}
                  \item 风险转移的时间相同。越过船舷（2000） \yjt 装运上船。注意“完成”交货和货交的区别。
                        \begin{enumerate}
                          \item \al{自收到、交给}。
                        \end{enumerate}
                  \item 均属于装运合同。装运港交货。
                  \item FOB(Free On Board) + 装运港。不包运、险。成本。卖方保险资料 + 通知投保 ⇄ 买方通知装运。
                  \item CFR(Cost and Freight) + 目的港。只包运不包险。成本 + 运费。卖方通知投保。
                  \item CIF(Cost Insurance and Freight) + 目的港。包运包险。成本 + 运费 + 保费（平安险）。
                  \item \sjt 只能用于船运 + FAS(Free Alongside Ship, 船边交货)（和 FOB 区别就是装运港装不装货）。
                  \item 关联性的 3 个。FCA(Free Carrier)、CPT(Carriage Paid to)、CIP(Carriage and Insurance Paid to)。
                        \begin{enumerate}
                          \item 适用所有运输方式。装运地交货。货交第一承运人时转移风险。
                          \item 其余各方面 FOB/FCA、CFR/CPT、CIF/CIP 一一对应。
                          \item FCA 的特殊规定。买卖双方可约定买方指示其承运人在货物装船后向卖方签发已装船提单。
                          \item CIP。除非另有约定，卖方需要投保的险别提高到一切险。
                        \end{enumerate}
                \end{enumerate}
          \item DAP(Delivered At Place) 和 DPU(Delivered at Place Unloaded)。所有 D 打头的卖方交货义务。
                \begin{enumerate}
                  \item 后跟目的地。卖方需要运到目的地（不限于运输终点、终端）。
                  \item 区别在运到后是否卸货。风险转移也是，一个运到后，一个卸货后。
                \end{enumerate}
        \end{enumerate}
\end{enumerate}
